\section{Regarding NHST}\label{app:nhst}
Throughout this study the focus was on three methods that use posterior information
to not only reject the null hypothesis but also accept it, or declare an inconclusive
result.

In contrast, null hypothesis significance testing (NHST) cannot accept the null
hypothesis. A p-value measures the probability of observing data at least as extreme
as those obtained \emph{assuming the null is true}, $P(\mathrm{data}\mid H_0)$,
not the probability that the null is true given the data,
$P(H_0\mid\mathrm{data})$
% [cohen1994] Cohen (1994), ``The earth is round (p < .05)'', American Psychologist
% 49(12):997--1003. The canonical demonstration that p-values cannot be read as
% P(H0|data). Directly supports the claim that NHST carries no null-acceptance
% information.
%
% [wasserstein2016] Wasserstein & Lazar (2016), ``The ASA Statement on p-Values:
% Context, Process, and Purpose'', The American Statistician 70(2):129--133.
% Official ASA position: ``A p-value does not provide a good measure of evidence
% regarding a model or hypothesis.'' Adds broad authoritative institutional backing.
\citep{cohen1994, wasserstein2016}.
In NHST, at each iteration the experiment stops and $\theta_{\rm null}$ is rejected
whenever the p-value falls below the false positive rate threshold $\alpha_{\rm FPR}$.
Since the decision rule is based solely on this single threshold, data that fail to
trigger rejection yield an inconclusive result rather than evidence in favour of the null.

Using the same fair coin setup ($\theta_{\rm true}=\theta_{\rm null}=0.5$) described
in the Methods section (\ref{sec:methods}), we apply NHST with $\alpha_{\rm FPR}=0.05$
as the stopping and decision criterion. A reliable sequential method should maintain a
low rejection rate throughout when the null hypothesis is true: as demonstrated in
Section~\ref{sec:fair_coin}, HDI+ROPE's rejection rate quickly asymptotes near 6\% and
both precision-based methods maintain a zero rejection rate.

Figure~\ref{fig:fair_decisions_nhst} shows the opposite for NHST: the rejection rate
(red solid) grows steadily over time, while the inconclusive proportion (gray dashed)
decreases. To make this trend visible, the simulation is extended to
$N_{\rm max}=30{,}000$ iterations, well beyond the $N_{\rm max}=1{,}500$ used in
the main analysis, since false rejections accumulate slowly and the growing trend
would not be apparent at shorter horizons.

\begin{figure}[h!]
  \centering
  \includegraphics[width=1\textwidth]{fair_experiment_decision_rates_nhst.png}
  \caption{Cumulative decision rates across iterations for $M=50$ fair coin experiments
  ($\theta_{\rm true}=\theta_{\rm null}=0.5$, $\alpha_{\rm FPR}=0.05$),
  analogous to Figure~\ref{fig:fair_decisions} but using NHST as the stopping and
  decision criterion, extended to $N_{\rm max}=30{,}000$ iterations.
  At each iteration $N$, the plotted proportions count all experiments that reached
  a decision at or before $N$; their sum is 100\% at every point.
  \textit{Line styles}:
  red solid: reject $\theta_{\rm null}$;
  gray dashed: inconclusive.
  The smaller experiment count ($M=50$ vs.\ $2{,}000$) reflects the higher
  computational cost of evaluating p-values relative to HDI-based criteria,
  which also accounts for the stepped appearance of the curves.}
  \label{fig:fair_decisions_nhst}
\end{figure}

As pointed out by \citet{kruschke2015doing}, as evidence accumulates under NHST
the probability of observing a sample extreme enough to trigger rejection grows without
bound --- even when the null hypothesis is true. To paraphrase: ``with infinite
peeking, the false positive rate approaches 100\%.'' This is because the stopping
rule is tied directly to the decision rule with no reference to an effect size: the
more data collected, the more likely an extreme sample will appear. Continuously
monitoring the data and applying a fixed p-value threshold therefore leads to an
inflated Type~I (false positive) error rate, a manifestation of the multiple-testing
problem.

\section{Regarding Bayes Factors}\label{app:app_bayes_factors}

The Bayes Factors is another popular method that uses the posterior to reject or accept
the null hyothesis, but as opposed to the HDI+ROPE method and the precision based methods
does not involve regards to the effect size, the ROPE.
The Bayes Factor is a ratio of the likelihood of the data under the null hypothesis to
the likelihood of the data under the alternative hypothesis,
and it can be used to quantify evidence for or against the null hypothesis.
However, it does not provide a direct measure of precision or a way to set a precision goal,
which is a key feature of the methods discussed in this paper.


TODO: continue from here

TODO: replicate the Bayes Factor of Figure 18.1 of https://ethanweed.github.io/pythonbook/06.01-bayes.html
as well as that in 13.6 of \cite{kruschke2015doing}.